\chapter{Suites monotones, points fixes et suites adjacentes}
\section*{Exercice 11.1 -- Suite racine}
[Méthode] Étudier $u_0=1$, $u_{n+1}=\sqrt{3+u_n}$.

\section*{Exercice 11.2 -- Erreur contractante}
[Base] Si $|u_n-L|\le8(1/2)^n$, donner un rang garantissant une erreur au plus $1/100$.

\section*{Exercice 11.3 -- Point fixe trompeur}
[Critique] Pour $u_0=1$ et $u_{n+1}=-u_n$, résoudre l'équation du point fixe et comparer au comportement réel.

\section*{Exercice 11.4 -- Adjacentes}
[Preuve] Pour $a_n=1-1/(n+1)$, $b_n=1+1/(n+1)$, vérifier l'adjacence et encadrer leur limite.

\section*{Exercice 11.5 -- Héron pour 3}
[Méthode] Partant de $u_0=2$, calculer deux itérations de Héron pour $\sqrt3$ et expliquer la convergence.

\section*{Exercice 11.6 -- Limite ou majorant}
[Critique] Une suite croissante est majorée par 100. Peut-on conclure qu'elle tend vers 100 ?
